\documentclass[11pt]{article}
\usepackage[margin=1in]{geometry}
\usepackage{amsmath,amssymb,amsthm}
\newtheorem{theorem}{Theorem}
\title{A disproof of Conjecture 00000000752}
\author{earthking11}
\date{October 6, 2026}
\begin{document}
\maketitle
\section*{Counterexample}
Take $p=3$ and the full one-sided shift on two symbols.  A presentation by a
one-vertex graph has adjacency matrix $A=[2]$.  Its characteristic polynomial
is $X-2$.

\begin{theorem}
The shift has entropy $\log 2$, whereas the conjectural p-adic-root formula
gives $0$.
\end{theorem}
\begin{proof}
Every position of a word can independently be either symbol, so the number of
admissible length-$n$ words is $2^n$.  Consequently
\[
 h_{\rm top}=\lim_{n\to\infty}\frac1n\log(2^n)=\log 2>0.
\]
The only characteristic root is $2$.  It is a $3$-adic unit, hence
$|2|_3=1$.  The formula in the conjecture therefore returns $\log|2|_3=
\log1=0$.  These values differ.
\end{proof}

The issue is structural: real topological entropy is governed by the ordinary
spectral radius, not by taking a p-adic norm of the characteristic roots.

\section*{Lean 4 certificate}
The project in \texttt{lean4/} defines the full-shift word count as the finite
product $k^n$, computes $2^n$, certifies that $2$ is a $3$-adic unit of norm
$1$, and proves the predicted equality impossible at $n=1$.  No unproved
placeholder or additional axiom is used.
\end{document}
